\documentclass[10pt,a4paper]{article}
\usepackage[margin=1.65cm]{geometry}
\usepackage{amsmath,amssymb,amsthm,mathtools}
\usepackage{graphicx}
\usepackage[hidelinks]{hyperref}
\usepackage{microtype}
\usepackage{xcolor}

\newtheorem{theorem}{Theorem}
\newtheorem{lemma}[theorem]{Lemma}
\newtheorem{proposition}[theorem]{Proposition}
\newtheorem{remark}[theorem]{Remark}
\newcommand{\M}{M}
\newcommand{\R}{\mathbb R}
\newcommand{\e}{\mathrm e}
\newcommand{\Id}{I}
\newcommand{\tr}{\operatorname{tr}}
\setlength{\parindent}{0pt}
\setlength{\parskip}{0.42em}

\title{Explicit bidirectional separation of product and additive integrability\\
\large A full answer to the first open problem in arXiv:1410.1845}
\author{Autonomous mathematical discovery run \texttt{fa\_banach\_001}}
\date{17 August 2026}

\begin{document}
\maketitle

\begin{abstract}
We construct step mappings with well-ordered steps which separate Kurzweil
product integrability and Henstock--Kurzweil integrability in both directions.
The construction uses exact five-factor commutator blocks in \(M_3(\R)\):
their exponential product is the identity while their additive sum has a
harmonic nilpotent drift.  Removing the correcting fifth factor reverses the
separation.  A principal-log refinement gives the product-integrable but
non-Henstock--Kurzweil example already in \(M_2(\R)\).  We also isolate a
missing local-log condition in the sequence-level equivalence asserted in the
source, and verify the original step-integrability statements directly.
\end{abstract}

\section{Source question and answer}

The source is S. Heikkil\"a and A. Slav\'ik, \emph{On summability,
multipliability, product integrability, and parallel translation},
arXiv:1410.1845; J. Math. Anal. Appl. 433 (2016), 887--934.  Section 9,
PDF page 44, asks:

\begin{quote}
Is there a step mapping with well-ordered steps which is Kurzweil product
integrable but not Henstock--Kurzweil integrable (or vice versa)?
\end{quote}

\begin{center}
\includegraphics[width=0.94\textwidth]{figures/open_problem_crop.png}
\end{center}

The answer is yes in both directions, with completely explicit real matrix
algebras.  We use the source convention that ordered products accumulate new
factors on the left.

\section{Proof intuition}

Noncommutativity lets a short block cancel additively while retaining a
second-order multiplicative drift.  For nilpotent matrix units this drift is
exact: a four-factor commutator equals \(\Id-t^2W\), with no higher-order
remainder.  Repeating it at \(t=n^{-1/2}\) produces harmonic product drift
but zero additive drift.  Appending the inverse drift factor makes every
block product the identity while leaving a harmonic additive sum.  Embedding
the resulting sequences as shrinking steps transfers these two elementary
block calculations to the two requested integrability separations.

\section{The sequence-to-step device}

Let \((x_j)_{j\geq1}\) be a sequence in a finite-dimensional unital Banach
algebra \(E\), and put
\[
 \alpha_j=1-2^{1-j},\qquad
 \Lambda=\{\alpha_j:j\geq1\}\cup\{1\}.
\]
Then \(\Lambda\subset[0,1]\) is well ordered, its only limit element is \(1\),
and \(\alpha_{j+1}-\alpha_j=2^{-j}\).  Define the step mapping
\begin{equation}\label{eq:embed}
 F_x(t)=2^j x_j\quad(t\in[\alpha_j,\alpha_{j+1})),
 \qquad F_x(1)=0.
\end{equation}
Thus the additive increment on the \(j\)-th step is exactly \(x_j\), and its
exponential increment is \(\e^{x_j}\).

\begin{lemma}\label{lem:device}
Assume \(x_j\to0\).
\begin{enumerate}
\item If \(P_N=\e^{x_N}\cdots\e^{x_1}\) converges to an invertible \(P\),
then \(F_x\) is strongly Kurzweil product integrable and its product integral
is \(P\).
\item The mapping \(F_x\) is Henstock--Kurzweil integrable exactly when
\(\sum_jx_j\) converges.
\item If \((P_N)\) does not converge, then \(F_x\) is not Kurzweil product
integrable.
\end{enumerate}
\end{lemma}

\begin{proof}
For (1), set \(P_0=\Id\) and, for
\(t\in[\alpha_j,\alpha_{j+1})\), define
\[
 W(t)=\exp\!\big((t-\alpha_j)2^jx_j\big)P_{j-1},
 \qquad W(1)=P.
\]
At every finite step boundary the two definitions agree.  Moreover,
\(P_{j-1}\to P\) and
\(\|(t-\alpha_j)2^jx_j\|\leq\|x_j\|\to0\); hence \(W(t)\to P\) as
\(t\uparrow1\).  Thus \(W\) is continuous and invertible, and
\(W'(t)=F_x(t)W(t)\) off the countable set of step boundaries.  The continuous
fundamental-matrix criterion in source Theorem 4.1 gives strong Kurzweil
product integrability.

Statement (2) is the standard well-ordered-step criterion used by the source
(its cited Proposition 3.1): the endpoint values of the additive primitive are
\(\sum_{j=1}^N x_j\).  In finite dimension this is also immediate from
continuity of the indefinite Henstock--Kurzweil integral and the converse
construction of the piecewise-linear primitive.  Finally, a product-integrable
mapping with countably many discontinuities has a continuous indefinite
product integral by source Theorem 4.1.  Its values at \(\alpha_{N+1}\) are
\(P_N\), proving (3).
\end{proof}

\section{Exact nilpotent commutator blocks}

Work first in \(E=M_3(\R)\), with the operator norm, and write
\[
 U=E_{12},\qquad V=E_{23},\qquad W=E_{13}.
\]
Then \(UV=W\), \(VU=0\), and every product involving \(W\) and one of
\(U,V,W\) is zero.  Since \(U^2=V^2=W^2=0\), direct multiplication gives
the exact identity
\begin{equation}\label{eq:comm}
 \e^{-tV}\e^{-tU}\e^{tV}\e^{tU}=\Id-t^2W,
 \qquad
 \e^{t^2W}(\Id-t^2W)=\Id.
\end{equation}

\begin{theorem}\label{thm:main}
There are two explicit \(M_3(\R)\)-valued step mappings on \([0,1]\), each
with the well-ordered step set \(\Lambda\) above, such that:
\begin{enumerate}
\item the first is strongly Kurzweil product integrable, with product integral
\(\Id\), but is not Henstock--Kurzweil integrable;
\item the second is Henstock--Kurzweil integrable, with integral \(0\), but is
not Kurzweil product integrable.
\end{enumerate}
\end{theorem}

\begin{proof}
Put \(t_n=n^{-1/2}\).  For the first mapping, concatenate the five-term
blocks
\begin{equation}\label{eq:five}
 t_nU,\quad t_nV,\quad -t_nU,\quad -t_nV,\quad t_n^2W
 \qquad(n=1,2,\ldots)
\end{equation}
into a sequence \((x_j)\), and use \eqref{eq:embed}.  Each block has additive
sum \(t_n^2W=W/n\).  Consequently, at the block endpoints the additive
partial sums are \(H_NW\), where \(H_N=\sum_{n=1}^N1/n\), and hence the series
does not converge.

In the ordered exponential product, however, \eqref{eq:comm} shows that every
five-term block has product
\[
 \e^{t_n^2W}\e^{-t_nV}\e^{-t_nU}\e^{t_nV}\e^{t_nU}=\Id.
\]
The partial products inside the \(n\)-th block tend uniformly to \(\Id\),
because all five increments tend to zero.  Thus the full sequence of partial
products converges to \(\Id\).  Lemma \ref{lem:device} proves (1).

For the second mapping, concatenate only the first four terms in
\eqref{eq:five}, obtaining a sequence \((y_j)\).  Every additive block sum is
zero, and the largest norm of an intra-block additive partial sum is
\(O(t_n)\).  Therefore \(\sum_jy_j=0\).  On the other hand, the exponential
product of block \(n\) is \(\Id-W/n\).  Since \(W^2=0\), the product through
the first \(N\) blocks is exactly
\[
 (\Id-W/N)\cdots(\Id-W)=\Id-H_NW,
\]
which diverges in norm.  Lemma \ref{lem:device} proves (2).
\end{proof}

\section{Dimension-two refinement}

The first separation can be realized in the \(2\times2\) algebra highlighted
by the source.  Let \(A=E_{12}\), \(B=E_{21}\), and
\[
 C(t)=\e^{-tB}\e^{-tA}\e^{tB}\e^{tA}
 =\begin{pmatrix}1-t^2&-t^3\\ t^3&1+t^2+t^4\end{pmatrix}.
\]
We have \(\det C(t)=1\) and \(\tr C(t)=2+t^4>2\), so \(C(t)\) has positive
eigenvalues and a real principal logarithm.  Set
\[
 s(t)=\operatorname{arcosh}(1+t^4/2),\qquad
 D(t)=C(t)-(1+t^4/2)\Id.
\]
Since \(D(t)^2=\sinh^2(s(t))\Id\),
\begin{equation}\label{eq:log}
 \log C(t)=\frac{s(t)}{\sinh s(t)}D(t).
\end{equation}

\begin{proposition}\label{prop:m2}
There is an \(M_2(\R)\)-valued step mapping with well-ordered steps which is
strongly Kurzweil product integrable but not Henstock--Kurzweil integrable.
\end{proposition}

\begin{proof}
For \(t_n=n^{-1/2}\), concatenate the blocks
\[
 t_nA,\quad t_nB,\quad -t_nA,\quad -t_nB,
 \quad -\log C(t_n).
\]
Their ordered exponential products are exactly \(C(t_n)^{-1}C(t_n)=\Id\),
and every increment tends to zero.  Hence the partial products converge to
\(\Id\).

The first four additive terms cancel.  Apply the linear functional
\(\ell(X)=(X_{11}-X_{22})/2\).  Formula \eqref{eq:log} gives
\[
 \ell(-\log C(t))=
 \frac{s(t)}{\sinh s(t)}\left(t^2+\frac{t^4}{2}\right).
\]
The ratio \(s/\sinh s\) tends to \(1\) as \(s\to0\), so for all sufficiently
large \(n\), the displayed quantity is at least \(1/(2n)\).  The additive
series therefore diverges.  Apply Lemma \ref{lem:device}.
\end{proof}

\section{A source-equivalence caveat}

The source calls its sequence formulation equivalent to the step-mapping
question.  The proof of the converse in its Theorem 4.2 uses
\(\exp(x_\alpha)\to\Id\) to infer
\(x_\alpha=\log(\exp x_\alpha)\to0\).  This requires eventual membership in a
fixed logarithm chart and is false in general.  For
\[
 J=\begin{pmatrix}0&-1\\1&0\end{pmatrix},\qquad x_j=2\pi J,
\]
all exponentials equal \(\Id\), although \(x_j\not\to0\).  The corresponding
step mapping from \eqref{eq:embed} is not product integrable: its would-be
fundamental matrix executes a full rotation on every interval and has no limit
at \(1\).

This caveat does not affect Theorem \ref{thm:main} or Proposition
\ref{prop:m2}.  All their increments tend to zero, and product integrability
was proved from an explicitly continuous fundamental matrix.

\section{Verification, scope, and novelty}

The exact identities in \eqref{eq:comm} and the algebra behind
\eqref{eq:log} were independently checked by the included SymPy script.  The
proof answers the first source question in full and, in \(M_3(\R)\), in both
directions.  It does not address the three other open bullets on PDF page 44.

A bounded search of the run indexes and of exact-question, title,
product-integral, multipliability, matrix-product, correction, and erratum
phrases found no later answer or matching construction through 17 August 2026.
The source has few indexed citations, and a 2016 chapter by the same authors
was the closest follow-up.  Novelty is therefore plausible but should be
checked by a specialist.  Human review should focus on the source's precise
strong/ordinary conventions and on the use of its Theorem 4.1 in Lemma
\ref{lem:device}.

\begin{thebibliography}{9}
\bibitem{HS}
S. Heikkil\"a and A. Slav\'ik,
\emph{On summability, multipliability, product integrability, and parallel
translation}, arXiv:1410.1845; J. Math. Anal. Appl. 433 (2016), 887--934.

\bibitem{S}
A. Slav\'ik,
\emph{Kurzweil and McShane product integration in Banach algebras},
J. Math. Anal. Appl. 424 (2015), 748--773.
\end{thebibliography}

\end{document}
